% ============================================================
% METODOLOGÍA — TFG
% Índice Sintético de Turismo Sostenible por CCAA (2016–2024)
% ============================================================

\chapter{Metodología}

\section{Datos y construcción del panel}

Los datos provienen de la plataforma estadística EVASTUR, desarrollada por Dephimática, que centraliza información turística regional procedente de diversas fuentes oficiales. El período analizado abarca de 2016 a 2024. El conjunto inicial contiene 58 indicadores organizados en tres dimensiones: económica (ECO), ambiental (ENV) y social (SOC), referidos a 19 comunidades autónomas españolas más un agregado nacional (ESP).

Todo el pipeline analítico se implementa en R mediante el paquete \texttt{COINr} \citep{becker2022coinr}, que estructura los datos como un objeto \emph{purse} —una colección de \emph{coins} anuales—. Cada \emph{coin} contiene una matriz de datos de dimensión 20 filas (19 CCAA + España) $\times$ número de indicadores por año.

\section{Selección y depuración de indicadores}

\subsection{Criterio 1: Varianza nula}

Se identifican indicadores cuyo rango interanual es nulo en al menos un año, es decir, aquellos para los que $\min = \max$. Este fenómeno causa división por cero en la normalización min-max:

\begin{equation}
    x'_i = 1 + \frac{(x_i - x_{\min})(b - a)}{x_{\max} - x_{\min}}
    \quad \Rightarrow \quad \text{indefinido si } x_{\max} = x_{\min}
\end{equation}

El análisis sobre 9 años identifica 11 indicadores con este problema en al menos un año. Cuando se excluye 2020 (año de disponibilidad limitada), el número se reduce a 5: IAMB0003, IAMB0005, IECO0012, IECO0013 e ISOC0004.

\subsection{Criterio 2: Disponibilidad de datos}

Se calcula el porcentaje de observaciones no ausentes (NA) por indicador:

\begin{equation}
    \text{Disponibilidad}_j = \frac{\text{# observaciones} \, \text{con valor}}{\text{# total de celdas posibles}} \times 100\,\%
\end{equation}

donde el denominador es 19 CCAA $\times$ 9 años = 171. Se descartan indicadores con disponibilidad $< 20\,\%$, criterio que elimina 6 indicadores adicionales: IAMB0004, IAMB0003, IAMB0005, IECO0012, IECO0013 e ISOC0004.

\textbf{Resultado}: panel depurado con $\mathbf{52 \text{ indicadores}}$ activos.

\section{Imputación de valores ausentes mediante MICE}

\subsection{Justificación del método}

El panel contiene valores ausentes de naturaleza heterogénea. La imputación simple por mediana o eliminación por lista completa subestima la variabilidad cuando hay heterogeneidad. Se emplea \textbf{imputación múltiple por ecuaciones encadenadas} (MICE, del paquete \texttt{mice} de R).

\subsection{Especificación técnica}

MICE es un algoritmo iterativo Gibbs sampling. Se emplea el algoritmo \textbf{Predictive Mean Matching} (PMM) con parámetros:
\begin{itemize}
    \item $m = 1$ (una única cadena imputada)
    \item \texttt{method = "pmm"} (Predictive Mean Matching busca 5 vecinos más próximos)
    \item \texttt{maxit = 5} (iteraciones)
    \item \texttt{seed = 42} (reproducibilidad)
\end{itemize}

PMM asegura que los valores imputados son reales (rango empírico preservado) y que la distribución marginal se mantiene sin introducir artifacts de normalidad forzada.

\textbf{Resultado}: reducción de 3.065 NA (32,75\%) a 118 NA (1,26\%), con residuales únicamente en IAMB0007.

\section{Tratamiento de valores atípicos}

\subsection{Diagnóstico mediante asimetría y curtosis}

Se emplean dos estadísticos:

\begin{equation}
    \text{Skewness} = \frac{E[(X - \mu)^3]}{\sigma^3}, \quad \text{Kurtosis} = \frac{E[(X - \mu)^4]}{\sigma^4}
\end{equation}

con criterios del manual OCDE:
\begin{itemize}
    \item $|\text{Skew}| > 2$ indica cola larga significativa
    \item $\text{Kurt} > 3.5$ indica distribución leptocúrtica
\end{itemize}

\subsection{Winsorización}

Se realiza mediante la función \texttt{Treat()} de COINr. La winsorización reemplaza valores extremos por el umbral permitido:

\begin{equation}
    x_i^{\text{wins}} = \begin{cases}
        p_L & \text{si } x_i < p_L \\
        x_i & \text{si } p_L \leq x_i \leq p_U \\
        p_U & \text{si } x_i > p_U
    \end{cases}
\end{equation}

Parámetros: \texttt{winmax = 5}, \texttt{skew\_thresh = 2}, \texttt{kurt\_thresh = 3.5}, \texttt{force\_win = TRUE}. El tratamiento se aplica año a año (por \emph{coin}).

\section{Normalización: Min-Max en rango $[1, 100]$}

\begin{equation}
    x'_i = 1 + 99 \cdot \frac{x_i - x_{\min}}{x_{\max} - x_{\min}}
\end{equation}

donde $x_{\min}$ y $x_{\max}$ se calculan sobre el panel completo. Se elige rango [1, 100] porque la media geométrica requiere valores $> 0$ (logaritmo definido).

\section{Análisis de componentes principales (PCA) para ponderación}

\subsection{Estructura jerárquica en dos niveles}

\textbf{Nivel 1}: PCA dentro de cada dimensión para obtener pesos de indicadores. Para la dimensión $d$, los pesos normalizados se basan en los loadings absolutos de PC1:

\begin{equation}
    w_{j}^{(1)} = \frac{|\lambda_{j,1}|}{\sum_{i=1}^{k_d} |\lambda_{i,1}|}
\end{equation}

\textbf{Nivel 2}: PCA sobre los scores dimensionales:

\begin{equation}
    S_{d,i} = \sum_{j \in d} w_j^{(1)} \cdot x'_{ij}
\end{equation}

Los pesos de Nivel 2 se derivan análogamente.

\section{Métodos de agregación: cinco casos}

Se comparan cinco combinaciones metodológicas:

\subsection{Caso 1: Min-Max + PCA (PC1) + Media Aritmética}

\begin{equation}
    I_u^{(1)} = \sum_{j=1}^{J} w_j \cdot x'_{uj}
\end{equation}

donde $w_j$ incluye pesos de ambos niveles. Propiedades: compensabilidad total, aditividad lineal.

\subsection{Caso 2: Min-Max + PCA (PC1) + Media Geométrica}

\begin{equation}
    I_u^{(2)} = \exp\left(\frac{\sum_{j=1}^{J} w_j \ln x'_{uj}}{\sum_{j=1}^{J} w_j}\right)
\end{equation}

Propiedades: no-compensabilidad parcial, penaliza desequilibrios, rango $[1, 100]$.

\subsection{Caso 3: Z-Score + PCA (PC1) + Media Aritmética}

Estandarización:

\begin{equation}
    x'_i = \frac{x_i - \bar{x}}{\sigma_x}
\end{equation}

Agregación mediante media aritmética. Equivalente a Caso 1 en términos de ranking bajo ponderación PCA lineal.

\subsection{Caso 4: Min-Max + PCA (ALL componentes) + Media Geométrica ponderada}

Este caso utiliza **todas** las componentes principales (no solo PC1). Para cada dimensión:

\begin{equation}
    \text{Var. explicada por PC}_i = \frac{\sdev_i^2}{\sum_k \sdev_k^2}
\end{equation}

Los scores de cada PC se ponderan por su varianza explicada y se combinan mediante media geométrica. Posteriormente, los índices dimensionales se agregan mediante:

\begin{equation}
    I_u^{(4)} = (ECO_u^{\alpha} \cdot SOC_u^{\beta} \cdot ENV_u^{\delta})^{1/(\alpha+\beta+\delta)}
\end{equation}

donde $\alpha$, $\beta$, $\delta$ son las varianzas explicadas por PC1 de cada dimensión. Este enfoque captura toda la variablidad de datos, no solo la dirección dominante.

\subsection{Caso 5: PCA/FA rotado (método OCDE) + Media Aritmética}

Se retienen $k$ factores con Eigenvalue $> 1$ y se aplica rotación Varimax. Los loadings se concentran: cada indicador carga fuertemente en un único factor. Pesos se calculan como:

\begin{equation}
    w_j^{\text{OCDE}} = \frac{\text{loading}_j^2}{\sum_{i \in \text{factor}} \text{loading}_i^2} \times \frac{\text{Var. explicada por factor}}{\sum_f \text{Var. explicada}}
\end{equation}

Agregación mediante media aritmética. Este enfoque sigue explícitamente el manual OCDE de construcción de indicadores compuestos.

\end{document}
